%------------------------------------------------------------------------------%

\usepackage{apresentacao}
\usepackage{mathconfig}
\usepackage{tabto}

\makeatletter
\graphicspath  {{./figs/}}
\def\input@path{{./figs/}}
\makeatother

\newcommand{\fmedio}{f_{\text{médio}}}

%------------------------------------------------------------------------------%
\title {Integrais Impróprias}
\author{\large Luis Alberto D'Afonseca}
\date  {Integrais e Séries}

%------------------------------------------------------------------------------%
\begin{document}

%------------------------------------------------------------------------------%
\begin{frame}
    \titlepage
\end{frame}

%------------------------------------------------------------------------------%
\section{Integrais Impróprias}

%------------------------------------------------------------------------------%
\begin{frame}{Integrais Impróprias}
\begin{enumerate}[<+->]
  \setlength\itemsep{1em}
  \item Intervalo de integração infinito
  \item Descontinuidade no extremo do intervalo
\end{enumerate}
\end{frame}

%------------------------------------------------------------------------------%
\begin{frame}{Integrais Impróprias}
\begin{enumerate}[<+->]
  \setlength\itemsep{1em}
  \item Envolve o calculo de um limite
  \item Se o limite existe a integral é \emph{convergente}
  \item Se algum limite não existir a integral é \emph{divergente}
\end{enumerate}
\end{frame}

%------------------------------------------------------------------------------%
\section{Integrais Impróprias -- Tipo 1}

%------------------------------------------------------------------------------%
\begin{frame}{Integrais Impróprias -- Tipo 1}
  $f$ função contínua em intervalo ilimitado,
  \emph{\([a, \infty)\)},
  \emph{\((-\infty, b]\)},
  \emph{\((-\infty, \infty)\)}
\end{frame}

%------------------------------------------------------------------------------%
\begin{frame}{Integrais Impróprias -- Tipo 1}
  $f$ é função contínua em \([a, \infty)\)
  \large
  \[
    \int_a^{{\color<2>{blue} \infty}} f(x) \, dx
    \;=\;  {{\color<2>{blue} \lim_{t\to\infty}}}
    \int_a^{{\color<2>{blue} t}} f(x)  \, dx
  \]
\end{frame}

%------------------------------------------------------------------------------%
\begin{frame}{Integrais Impróprias -- Tipo 1}
  $f$ é função contínua em \((-\infty, b]\)
  \large
  \[
    \int_{{\color<2>{blue} -\infty}}^b f(x) \, dx
    \;=\; {\color<2>{blue} \lim_{t\to-\infty}}
    \int_{{\color<2>{blue} t}}^{b} f(x) \,dx
  \]
\end{frame}


%------------------------------------------------------------------------------%
\begin{frame}{Integrais Impróprias -- Tipo 1}

  $f$ é função contínua em \((-\infty, \infty)\)

  {\large
  \[
    \int_{{\color<3>{blue}  -\infty}}
        ^{{\color<4>{magenta}\infty}} f(x) \, dx
    \;=\; {\color<3>{blue}   \lim_{r\to-\infty}}
    \int_{{\color<3>{blue}{r}}}
        ^{{\color<2>{red} c}} f(x) \, dx
    \;+\; {\color<4>{magenta}\lim_{s \to \infty}}
    \int_{{\color<2>{red} c}}
        ^{{\color<4>{magenta}s}} f(x) \, dx
  \]}
  para algum {\color<2>{red} $c$} $\in\R$

\end{frame}

%------------------------------------------------------------------------------%
\section{Integrais Impróprias -- Tipo 2}

%------------------------------------------------------------------------------%
\begin{frame}{Integrais Impróprias -- Tipo 1}
  $f$ função \emph{contínua no interior} do intervalo,
  mas tende ao \emph{infinito} em algum ponto
\end{frame}

%------------------------------------------------------------------------------%
\begin{frame}{Integrais Impróprias -- Tipo 2}
  $f$ é contínua em {\color{blue} $[a, b)$} e descontínua em $b$,
  {
  \large
  \[
    \int_{a}^{{\color<2>{blue} b}} f(x)\,dx
    \;=\;     {\color<2>{blue} \lim_{t\to b^{-}}}
    \int_{a}^{{\color<2>{blue} t}} f(x)\,dx
  \]}
  Atenção ao limite lateral
\end{frame}

%------------------------------------------------------------------------------%
\begin{frame}{Integrais Impróprias -- Tipo 2}
  $f$ é contínua em {\color{blue} $(a, b]$} e descontínua em $a$,
  {
  \large
  \[
    \int_{{\color<2>{blue} a}}^b f(x)\,dx
    \;=\; {\color<2>{blue} \lim_{t\to a^+}}
    \int_{{\color<2>{blue} t}}^b f(x)\,dx
  \]}
  Atenção ao limite lateral
\end{frame}

%------------------------------------------------------------------------------%
\begin{frame}{Integrais Impróprias -- Tipo 2}
  $f$ é uma função definida em $[a, b]$, contínua exceto em um ponto
  {\color{blue} $c \in (a, b)$}
  {
  \large
  \begin{align*}
      \int_a^b f(x) \, dx
    & \;=\; \int^{{\color<2>{blue} c}}_a f(x) \, dx
      \;+\; \int_{{\color<2>{blue} c}}^b f(x) \, dx \\[3mm]
    & \;=\; {\color<2>{blue} \lim_{r\to c^-}} \int^{{\color<2>{blue} r}}_a f(x) \, dx
      \;+\; {\color<2>{blue} \lim_{s\to c^+}} \int_{{\color<2>{blue} s}}^b f(x) \, dx
  \end{align*}

  }
\end{frame}

%------------------------------------------------------------------------------%
\section{Exemplo}

%------------------------------------------------------------------------------%
\begin{frame}{Exemplo}
\onslide<+->{Calcular a integral \(\int_{-\infty}^0 2^x\, dx\)}
\begin{align*}
  \onslide<+->{\int_{{\color<2>{blue}-\infty}}^0 2^x\, dx
    & = {\color<2>{blue}\lim_{t\to-\infty}}
        \left[\;\int_{{\color<2>{blue} t}}^{0} 2^x\, dx \;\right]
    & & \text{definição integral imprópria}                       \\}
  \onslide<+->{
    & = \lim_{t\to-\infty}
        \left[\;{\color<3>{blue} \frac{2^x}{\ln(2)}} \,\evalat{t}{0} \;\right]
    & & \int a^u \,du = \frac{a^u}{\,\ln(a)\,} + C \\}
  \onslide<+->{
    & = \lim_{t\to-\infty} \left[\;\frac{1}{\ln(2)} \left( 2^0 - 2^t \right) \;\right] \\}
  \onslide<+->{
    & = \frac{1}{\ln(2)} \lim_{t\to-\infty} \left( 2^0 - 2^t \right) }
  \onslide<+->{
      \;=\; \frac{1}{\ln(2)}}
\end{align*}
\end{frame}

%------------------------------------------------------------------------------%
\begin{frame}
  \tableofcontents
\end{frame}

%------------------------------------------------------------------------------%
\end{document}
%------------------------------------------------------------------------------%
